\documentclass{article}
\usepackage[equations]{propositions}
\usepackage{setspace}

%% PROTOTYPE of a `conclusion' style: a rule across the text width, drawn
%% immediately above the list, for the conclusion of an argument whose
%% premises are in the list before it.
%%
%% Two requirements, neither of which a plain \hrule meets:
%%
%%   (i)  the rule adds no vertical space -- the first item must sit exactly
%%        where it would have sat without it;
%%   (ii) the rule sits midway between the last line above and the first line
%%        of this list.
%%
%% How it is done.  The wrapper-begin code runs in vertical mode, after the
%% package has deposited the seamless-join glue (see continue=true) and
%% immediately before \begin{list}.
%%
%% The rule cannot simply be added there.  The join works by depositing the
%% glue and then setting \topsep to match, so that \@trivlist's
%% \addvspace{\@topsepadd} finds \lastskip already equal to it and adds
%% nothing.  A box appended after that glue makes \lastskip zero, \addvspace
%% no longer sees its own value, and the join is laid down a second time --
%% \itemsep+\parsep of unwanted space.  So the rule is slipped *inside* the
%% glue: step back over it, drop the rule in, put the glue back.  \lastskip
%% is then what the list expects.
%%
%% Placement.  At that point \prevdepth is the depth of the line above and
%% the glue is \lastskip, so the next baseline falls \baselineskip past the
%% glue and the baseline-to-baseline gap is \baselineskip + <glue>.  The rule
%% belongs midway down that gap -- but midway between the *baselines* is not
%% midway to the eye, so the strut is used to find the white band instead (see
%% the code).  Having stepped back over
%% the glue, the insertion point is the last baseline itself -- measured, not
%% assumed: with the lowering set to zero the rule lands within half a point
%% of that baseline, and the residual does not track \prevdepth.  So the rule
%% is lowered by
%%
%%   0.5\baselineskip + 0.5<glue>
%%
%% The box has zero height, depth and width, \nointerlineskip stops TeX
%% adding glue above it, and \prevdepth is restored so the first line of the
%% list is spaced exactly as if the rule were not there.
%%
%% <glue> is \lastskip when the list is joining a previous one.  Standing on
%% its own after ordinary text there is no join glue, and the separation will
%% be the list's own \topsep+\partopsep instead.

%% The style and both commands now come from the package; this file only
%% exercises them.  The prototype that developed them, and the reasoning
%% behind the placement, is in the implementation section of propositions.dtx.

%% Markers, so the measuring script can find the lines either side.
\newcommand\PA{PREMISEA}
\newcommand\PB{PREMISEB}
\newcommand\CC{CONCLUSION}

\begin{document}

\section{The intended use}

An argument, with the conclusion ruled off from the premises:

\begin{prop}
  \item Everything that is \PA{} is physical.
  \item Everything that is mental is \PB.
\end{prop}
\begin{prop}[style=conclusion, continue=true]
  \item Everything that is mental is \CC.
\end{prop}

The rule must sit midway between the last premise and the conclusion, and the
conclusion must sit exactly where it would have without the rule.

\section{The same argument with no rule, for comparison}

\begin{prop}
  \item Everything that is \PA{} is physical.
  \item Everything that is mental is \PB.
\end{prop}
\begin{prop}[continue=true]
  \item Everything that is mental is \CC.
\end{prop}

The conclusion's baseline here and in the previous section must be the same
distance below its premise.

\section{Several premises, and a multi-line conclusion}

\begin{prop}
  \item Everything that is \PA{} is physical.
  \item Everything that is physical is spatiotemporal.
  \item Everything that is mental is \PB.
\end{prop}
\begin{prop}[style=conclusion, continue=true]
  \item Everything that is mental is \CC, and this item runs on long enough to
    take a second line, so that we can see the rule is placed relative to the
    first line of the item rather than to the item as a whole.
\end{prop}

\section{A premise list whose last line is short}

\begin{prop}
  \item Everything that is \PA.
\end{prop}
\begin{prop}[style=conclusion, continue=true]
  \item Therefore \CC.
\end{prop}

\section{Without \texttt{continue}, after ordinary text}

Some ordinary text, and then a conclusion list on its own.  Here there is no
premise list to join, so the rule falls in the list's own \texttt{topsep}
instead; it should still take no vertical space.

\begin{prop}[style=conclusion]
  \item A conclusion with nothing above it. \CC
\end{prop}

\section{Named items}

\begin{prop}
  \item[Physicalism] Everything is physical. \PA
\end{prop}
\begin{prop}[style=conclusion, continue=true]
  \item[Conclusion] So much the worse for dualism. \CC
\end{prop}

\section{The whole argument inside another item}

The rule spans \verb|\linewidth|, so within a nested list it should span the
nested text width, not the page's.

\begin{prop}
  \item An outer item that contains an argument of its own:
  \begin{prop}
    \item Everything that is \PA{} is physical.
    \item Everything that is mental is \PB.
  \end{prop}
  \begin{prop}[style=conclusion, continue=true]
    \item Everything that is mental is \CC.
  \end{prop}
  \item A further outer item, to show the argument has closed properly.
\end{prop}

\section{With \texttt{itemsep} changed}

The gap the rule has to bisect is the previous list's
\texttt{itemsep}+\texttt{parsep}, so it should follow whatever those are set
to.  Widely spaced:

\begin{prop}[itemsep=14pt]
  \item Everything that is \PA{} is physical.
  \item Everything that is mental is \PB.
\end{prop}
\begin{prop}[style=conclusion, continue=true, itemsep=14pt]
  \item Everything that is mental is \CC.
\end{prop}

Tightly spaced:

\begin{prop}[itemsep=0pt, parsep=0pt]
  \item Everything that is \PA{} is physical.
  \item Everything that is mental is \PB.
\end{prop}
\begin{prop}[style=conclusion, continue=true, itemsep=0pt, parsep=0pt]
  \item Everything that is mental is \CC.
\end{prop}

\section{Double spacing}

\textsf{setspace}'s \verb|\doublespacing| widens \verb|\baselineskip| and rescales
the list dimensions to match, so both terms of the calculation move together
and the rule should still land midway.

\begin{spacing}{2}
\begin{prop}
  \item Everything that is \PA{} is physical.
  \item Everything that is mental is \PB.
\end{prop}
\begin{prop}[style=conclusion, continue=true]
  \item Everything that is mental is \CC.
\end{prop}
\end{spacing}

One-and-a-half spacing, for a second point of comparison:

\begin{spacing}{1.5}
\begin{prop}
  \item Everything that is \PA{} is physical.
  \item Everything that is mental is \PB.
\end{prop}
\begin{prop}[style=conclusion, continue=true]
  \item Everything that is mental is \CC.
\end{prop}
\end{spacing}

\section{At the top of a page}

If the page breaks between the premises and their conclusion, the conclusion
list starts a page: there is no line above it, and \TeX\ has discarded the
glue that would have separated them.  There is nothing to centre between, so
the rule goes at the top and the clearance it would have had lower down is put
below it.  The gap between the rule and the text here should match the gap
below the rule in the sections above.

\begin{prop}
  \item A premise at the foot of the page. \PA
\end{prop}
\newpage
\begin{prop}[style=conclusion]
  \item A conclusion carried to the next page. \CC
\end{prop}

Text after, so the page is not empty.

\section{Rules above and below}

\verb|\propsetbelow| is the mirror of \verb|\propsetabove| and belongs in
\texttt{wrapper end}.  A style using both boxes a list in with rules rather
than a frame.  The two rules should sit the same distance from the text they
bracket.

\begingroup
\SetPropStyle{bracketed}{%
  wrapper begin = \propsetabove[0.2em][0.2em]{\smash{\rule{\linewidth}{0.4pt}}},
  wrapper end   = \propsetbelow[0.2em][0.2em]{\smash{\rule{\linewidth}{0.4pt}}}}

Some text before the bracketed list.

\begin{prop}[style=bracketed]
  \item Everything that is \PA{} is physical.
  \item Everything that is mental is \PB.
\end{prop}

Some text after the bracketed list, to show where the lower rule falls.
\endgroup

\section{Tuning: the optional argument}

The material is centred between the descenders of the line above and the tops
of the letters below.  Where something taller intrudes --- an opening
parenthesis, say, which rises above cap height --- the result can read a
little low.  The two optional arguments of \verb|\propsetabove| are the space to leave
above and below the material; either may be negative.  Equal values widen the
gap and keep the material centred in it; unequal values shift it.

\newcommand\showraise[3]{%
  \SetPropStyle{conclusion}{wrapper begin =
    \propsetabove[#1][#2]{\smash{\rule{\linewidth}{0.4pt}}}}%
  \begin{prop}
    \item Everything that is mental is PREMISEB.
  \end{prop}
  \begin{prop}[style=conclusion, continue=true]
    \item Everything that is mental is CONCLUSION \textup{(#3)}.
  \end{prop}}

\showraise{0.2em}{0.2em}{even, the default}
\showraise{0pt}{0pt}{no space either side}
\showraise{0.4em}{0.4em}{twice the space}
\showraise{0.4em}{0pt}{all the space above}
\SetPropStyle{conclusion}{wrapper begin =
  \propsetabove[0.2em][0.2em]{\smash{\rule{\linewidth}{0.4pt}}}}

\end{document}
